\documentclass[12pt]{article}
\usepackage{fullpage}
\begin{document}
	\section*{5 CONCLUSIONS}
	\subsection*{5.1 Introduction}
	This chapter presents the summary and the conclusions of the study in the subsections that follow.
	\subsection*{5.2 Summary}
	To identify the significant predictors that contribute to depression among cancer
	caregivers in a rural setting. To assess the performance of the multiple linear regression model in predicting depression among cancer caregivers. This involves evaluating the goodness of fit measures.To provide evidence-based recommendations for interventions and support systems
	aimed at addressing depression among cancer caregivers
	\subsection*{5.3 Conclusions}
	In conclusion, the age of patient and level of income are significant in determining the level of depression of cancer caregivers in that the level of depression decreases with increase in age of a patient as well as increase of level of income of a caregiver. The age of the cancer caregiver is insignificant.\\
	
	\noindent The factors of age of the patient and income of the caregiver account for 43\% and 77\% of the overall fitness, respectively. This suggests that income has a stronger correlation with the  level of depression than age. However, there may be other qualitative factors not considered in the analysis that contribute to the remaining percentage of fitness such as suppot groups and level of education.\\
	
	\noindent The outcomes of this study have practical implications for supporting caregivers in rural areas who experience depression. By understanding the factors influencing depression in this context, interventions can be designed to address these factors and provide effective support. This research contributes to the growing body of knowledge on mental health in caregiving situations and highlights the importance of considering the specific needs of caregivers in rural settings.\\
	
	\newpage
	\subsection*{5.4 Recommendation}
	Based on the findings of this research project, several evidence-based recommendations can be made to inform the development of targeted interventions for depression among cancer caregivers in rural settings:\\
	\begin{enumerate}
		\item We recommend fo more variables to be used in future research to get a more acccurate model i.e. qualitative variable support groups,level of education et al \\
		
		\item For easier avalability of cancer caregivers data in Kenya, Kenya National Bureau  of Statistics to conduct such research and make the data avalaible online for future research.\\
		
		\item Given the role being performed by cancer caregivers, the government through state departments should develop intervetions to provide tailored support to cancer caregivers for better quality life of the patients.\\  
		
		\item Provide financial assistance and resources: The level of income emerged as an important factor in the regression analysis. Interventions should address the financial challenges faced by caregivers in rural areas, such as limited job opportunities or transportation costs. Offering financial assistance programs or connecting caregivers with available resources can help alleviate financial stress and improve their mental well-being.\\
		
		
		
		
	\end{enumerate}
	
	
	
\end{document}